\documentclass{article}

% For anonymous main-track submission
\usepackage{neurips_2026}

% If you want a non-anonymous preprint instead, use:
% \usepackage[preprint]{neurips_2026}

% If you want camera-ready later, use:
% \usepackage[final]{neurips_2026}

\usepackage[utf8]{inputenc}
\usepackage[T1]{fontenc}
\usepackage{hyperref}
\usepackage{url}
\usepackage{booktabs}
\usepackage{amsfonts}
\usepackage{nicefrac}
\usepackage{microtype}
\usepackage{xcolor}
\usepackage{amsmath,amssymb,mathtools,amsthm}
\usepackage{algorithm}
\usepackage{algpseudocode}
\usepackage{graphicx} % needed for \resizebox

\title{Constructing Machine-Precision Neural Networks with Quasi-Interpolants}

% Anonymous submission
\author{%
  Anonymous Authors \\
  Paper under double-blind review
}

\begin{document}

\maketitle

\begin{abstract}
% Paste your abstract here
\end{abstract}

\section{Introduction}
% Paste your introduction here

\section{Related Work}
% Paste your related work here

% =========================
% SECTION 3
% =========================

\section{Constructing High-Precision MLPs with Quasi-Interpolants}
\label{sec:qi-construction}
Our goal is to build an interpolation scheme with arbitrarily good approximation that is also expressible in the primitives available to one-hidden-layer MLPs. Quasi-interpolants are a natural choice because they reconstruct from samples using translates of a single cardinal function, which will later construct directly with hidden units and activations. We first construct the quasi-interpolant and its cardinal function here, then realize it as an MLP (Section~3.2), derive the error-width tradeoff (Section~3.3), and summarize the implementation in Section~3.4.

\paragraph{Main result (informal).}
For any analytic target \(f\) on \([-1,1]\) and any target precision \(\varepsilon > 0\) above a kernel-dependent threshold, we construct a one-hidden-layer MLP with \(W = O(\log(1/\varepsilon))\) neurons and maximum weight magnitude \(O(\log(1/\varepsilon))\) that interpolates \(f\) to accuracy \(\varepsilon\) in \(L^\infty\). The construction is explicit (Algorithm~\ref{alg:qi-mlp}); for the \(\operatorname{sech}^2\) kernel, the viable parameter regime extends to fp64 machine precision (\(\varepsilon \sim 10^{-15}\)), with all weights remaining within floating-point range (Table~\ref{tab:comparison}).

\subsection{Constructing Fourier-Normalized Quasi-Interpolants}
\label{subsec:fourier-normalized}

\paragraph{Setup.} To begin our Quasi-interpolation construction from MLPs, we let \(\Omega = [-1,1]\), fix \(N \in \mathbb{N}\), and define
\begin{equation}
h := \frac{2}{N}, \qquad x_k := -1 + kh, \qquad K_\gamma(x) := \gamma K(\gamma x).
\label{eq:grid-kernel}
\end{equation}
A translation-invariant quasi-interpolant reconstructs \(f\) from uniform samples as
\begin{equation}
(Q_h f)(x) := \sum_{k \in \mathbb{Z}} f(x_k)\, L_h(x - x_k),
\label{eq:Qh}
\end{equation}
where the ideal \emph{cardinal function} \(L_h\) would satisfy the following two conditions
\begin{itemize}
    \item[\textbf{(i)}] \(\boldsymbol{L_h(jh) = \delta_{j,0}}\) for all \(j \in \mathbb{Z}\), ensuring exact interpolation at grid points.
    \item[\textbf{(ii)}] \(\boldsymbol{L_h(x) = \sum_j c_j\, K_\gamma(x - jh)}\), so the operator maps onto a one-hidden-layer MLP (Section~3.2).
\end{itemize}
Furthermore, stability between nodes is guaranteed by the kernel's exponential spatial decay (Assumption K4, Appendix A.1), which makes \(L_h\) a localized bump, the quasi-interpolation analog of avoiding the Runge phenomenon.

\paragraph{Fourier normalization.}
By Appendix~B.1, condition \textbf{(i)} is equivalent to the Fourier partition-of-unity condition
\begin{equation}
\sum_{m \in \mathbb{Z}} \widehat{L}_h\!\left(\omega + \frac{2\pi m}{h}\right) = h
\qquad \forall\, \omega.
\label{eq:partition-unity}
\end{equation}
And enforcing the structure found in condition \textbf{(ii)} lets us decouple $\widehat{L}_h(\omega)$ via shift invariance, giving us
\[
\widehat{L}_h(\omega)
=
\widehat{K}_\gamma(\omega)\sum_j c_j e^{-ijh\omega}
=
\widehat{K}_\gamma(\omega)\,\widehat{C}_h(\omega),
\qquad
\widehat{C}_h(\omega) := \sum_j c_j e^{-ijh\omega}.
\]
Since \(\widehat{C}_h\) is \(2\pi/h\)-periodic, plugging this into \eqref{eq:partition-unity} gives
\[
h
=
\widehat{C}_h(\omega)\sum_{m \in \mathbb{Z}} \widehat{K}_\gamma\!\left(\omega + \frac{2\pi m}{h}\right)
=
\widehat{C}_h(\omega)\,D_h(\omega),
\]
so the construction reduces to a single division in Fourier space:
\begin{equation}
\widehat{C}_h(\omega) = \frac{h}{D_h(\omega)},
\qquad
D_h(\omega) := \sum_{m \in \mathbb{Z}} \widehat{K}_\gamma\!\left(\omega + \frac{2\pi m}{h}\right).
\label{eq:fourier-character}
\end{equation}
Here \(\widehat{C}_h\) is the \emph{Fourier character} and \(D_h\) is the \emph{normalizer}, the periodized kernel spectrum. The construction requires \(D_h \neq 0\) everywhere, which we verify for \(K = \tfrac{1}{2}\operatorname{sech}^2\) via diagonal dominance (Appendix A.1, Proposition 4). Rescaling \(\theta = h\omega\) gives \(\widetilde{C}_\lambda(\theta) := \widehat{C}_h(\theta/h) = \sum_j c_j e^{-ij\theta}\), so the Fourier coefficients of \(\widetilde{C}_\lambda\) are exactly the same coefficients \(c_j\) appearing in the original kernel expansion of \(L_h\). In practice, these coefficients are computed using the DFT shown explicitly in Algorithm 1.

\paragraph{From infinite to finite.}
To implement the operator on \(\Omega\), we introduce two truncations. \emph{Stencil truncation}: restrict the cardinal function to \(|j| \leq K_c\) terms,
\begin{equation}
L_h^{(K_c)}(x) := \sum_{|j| \leq K_c} c_j\, K_\gamma(x - jh).
\label{eq:Lh-trunc}
\end{equation}
\emph{Halo truncation}: restrict the lattice sum to \(k \in \{-R, \ldots, N+R\}\), placing \(R\) nodes on each side of \(\Omega\) to stabilize boundary effects (Figure~2a). The finite operator is
\begin{equation}
(Q_{h,R}^{(K_c)} f)(x) := \sum_{k=-R}^{N+R} f(x_k)\, L_h^{(K_c)}(x - x_k),
\qquad x \in \Omega,
\label{eq:finite-op}
\end{equation}
which re-indexes into a \emph{kernel network} (Appendix~B.3):
\begin{equation}
(Q_{h,R}^{(K_c)} f)(x)
=
\sum_{m=-R}^{N+R} a[m]\, K_\gamma(x - x_m),
\qquad
a[m] := \sum_{|j| \leq K_c} c_j\, f(x_{m-j}).
\label{eq:kernel-network}
\end{equation}

\subsection{MLPs Can Implement Quasi-Interpolants}
\label{subsec:mlp-implementation}

The kernel network \eqref{eq:kernel-network} uses \(K_\gamma\) as its basis, but standard MLP activations \(\psi\) (e.g., \(\tanh\)) typically do not satisfy the kernel requirements directly. Instead, a low-order derivative often does: \(\operatorname{sech}^2 = \tanh'\). Assume
\begin{equation}
K(u) = \psi^{(r)}(u)
\label{eq:activation-kernel}
\end{equation}
for some smooth \(\psi\) and order \(r \in \mathbb{N}\), with \(K\) satisfying the kernel conditions (Assumptions K1--K4). By the chain rule, \(K_\gamma(x - x_m) = \gamma^{-(r-1)} \frac{d^r}{dx^r} \psi(\gamma(x - x_m))\), so \eqref{eq:kernel-network} becomes
\begin{equation}
(Q_{h,R}^{(K_c)} f)(x)
=
\frac{d^r}{dx^r}\!\left(
\sum_{m=-R}^{N+R} \frac{a[m]}{\gamma^{r-1}}\, \psi\big(\gamma(x - x_m)\big)
\right).
\label{eq:derivative-form}
\end{equation}
Define the one-hidden-layer MLP
\begin{equation}
g_{\mathrm{MLP}}(x)
:=
\sum_{m=-R}^{N+R} w[m]\, \psi\big(\gamma(x - x_m)\big),
\qquad
w[m] := \frac{a[m]}{\gamma^{r-1}}.
\label{eq:gmlp}
\end{equation}
Then
\begin{equation}
g_{\mathrm{MLP}}^{(r)}(x) = (Q_{h,R}^{(K_c)} f)(x).
\label{eq:gmlp-identity}
\end{equation}
Applying the construction to \(f^{(r)}\), the MLP recovers \(f\) up to an integration polynomial:
\begin{equation}
\widetilde{f}(x) = p_{r-1}(x) + g_{\mathrm{MLP}}(x),
\label{eq:f-tilde}
\end{equation}
where \(p_{r-1}\) is fixed by boundary conditions. For \(\tanh\) (\(r = 1\), \(K = \operatorname{sech}^2\)), all kernel conditions are verified (Appendix A.1) and the polynomial reduces to a bias. Extending the construction to other activations (e.g., GELU, Swish) via higher-order derivatives is possible in principle but requires separate verification of the kernel conditions for each case.

\paragraph{Expressivity is not the bottleneck.}
Equations \eqref{eq:gmlp}--\eqref{eq:f-tilde} show that one-hidden-layer tanh MLPs can exactly represent the quasi-interpolant construction. Whether machine-precision interpolation is achievable reduces to whether the correct weights can be found, an optimization question.

\subsection{Error Analysis and Width-Precision Scaling}
\label{subsec:error-analysis}

\paragraph{The dimensionless bandwidth.}
We now introduce the dimensionless bandwidth parameter
\[
\lambda := \gamma h.
\]
Since the kernel width is on the order of \(1/\gamma\) and the grid spacing is \(h\), the quantity \(1/\lambda = 1/(\gamma h)\) can be interpreted as the approximate number of grid points that lie inside one kernel width. The parameter \(\lambda\) controls a fundamental tradeoff. When \(\lambda\) is large (narrow kernels), the cardinal function is easy to construct but the kernel's spectral energy above the Nyquist rate \(\pi/h\) folds back into \(D_h\), producing \emph{aliasing error}. When \(\lambda\) is small (wide kernels), aliasing vanishes but building a cardinal function from broad overlapping bumps demands large coefficients with precise cancellations, making the construction \emph{ill-conditioned}, reflected in error-bound prefactors that diverge algebraically as \(\lambda \to 0\) (Theorem~1). The optimal \(\lambda\) balances these effects.

\paragraph{Error decomposition.}
The total error decomposes into four terms by the triangle inequality (full proof in Appendix A.1, Theorem 5).

\medskip
\noindent\textbf{Theorem 1 (Numerical error bound).}
\emph{Let \(\lambda = \gamma h\). Assume \(f\) is analytic in a strip of half-width \(\rho\) around \(\Omega\), and assume the normalizer \(D_h\) is bounded away from zero on a strip of half-width \(\sigma\). Then}
\begin{equation}
\resizebox{\columnwidth}{!}{$\displaystyle
\|f - Q_{h,R}^{(K_c)} f\|_{L^\infty(\Omega)}
\le
\underbrace{C_1(\lambda)\, e^{-c_1/h}}_{\text{resolution}}
+
\underbrace{C_2\, e^{-c_2/\lambda^p}}_{\text{aliasing}}
+
\underbrace{C_3(\lambda)\, e^{-c_3 \lambda R}}_{\text{halo}}
+
\underbrace{C_4(\lambda)\, e^{-c_4 \lambda K_c}}_{\text{stencil}}.
$}
\label{eq:error-bound}
\end{equation}
\emph{The rates \(c_1, \ldots, c_4\) depend on the target's analyticity \(\rho\), the kernel's decay rate \(a_K\), and the normalizer's strip margin \(\sigma\). The prefactors \(C_1, C_3, C_4\) capture the stability cost of Fourier normalization and diverge algebraically as \(\lambda \to 0\) (explicit formulas in Appendix A.1, Eqs.~29--31). \(C_2\) is independent of \(\lambda\). For \(K = \operatorname{sech}^2\), \(p = 1\).}

Each term is exponential in a parameter we control. The resolution term decays as \(e^{-c_1 N/2}\) with the number of grid points; the aliasing term decays as \(e^{-c_2/\lambda^p}\), vanishing as the kernel widens relative to the grid (small \(\lambda\)); and the halo/stencil terms decay geometrically in \(R\) and \(K_c\), respectively, at rates coupled through \(\lambda\).

\paragraph{Width-precision scaling.}
To achieve a target tolerance \(\varepsilon\), we choose \(\lambda\) in the \emph{viable regime}: small enough that \(C_2 e^{-c_2/\lambda^p} \leq \varepsilon\) (bounding aliasing), yet large enough that the prefactors \(C_i(\lambda)\) remain moderate (bounding ill-conditioning). For a fixed precision target such as fp64, this viable regime is a computable interval \([\lambda_{\min}, \lambda_{\max}]\) depending on \(\varepsilon\), the kernel, and the analyticity width (Appendix~C.1). With any \(\lambda^*\) in this interval, the remaining three terms each require their control parameter to scale as \(O(\log(1/\varepsilon))\):
\[
N = O\!\left(\frac{1}{c_1}\log\frac{1}{\varepsilon}\right),
\qquad
R = O\!\left(\frac{1}{c_3\lambda^*}\log\frac{1}{\varepsilon}\right),
\qquad
K_c = O\!\left(\frac{1}{c_4\lambda^*}\log\frac{1}{\varepsilon}\right).
\]
Since the MLP has \(W = N + 2R + 1\) neurons, the total width scales as \(W = O(\log(1/\varepsilon))\). Equivalently,

\medskip
\noindent\textbf{Corollary 1 (Exponential convergence in width).}
\emph{Fix a target precision \(\varepsilon\) in the viable regime (Appendix~C.1) and choose \(\lambda^*\) accordingly. There exist \(A, \alpha > 0\) such that the constructed MLP with \(W\) neurons satisfies \(\|f - \widetilde{f}_W\|_{L^\infty} \leq A\, e^{-\alpha W}\). (Proof in Appendix~C.2.)}

This matches the rate of classical polynomial interpolation for analytic targets (Trefethen, 2019), realized by a one-hidden-layer MLP with explicit weights.

\paragraph{Weight scaling.}
Maintaining \(\lambda^*\) constant as \(h \to 0\) requires \(\gamma = \lambda^*/h = \lambda^* N/2\): the inner-layer weights must grow proportionally with width. The maximum weight magnitude is \(O(W) = O(\log(1/\varepsilon))\), keeping all parameters within fp64 range even at machine precision (Appendix~C.3). This is what makes the construction numerically realizable, in contrast to prior results where weights scale as \(\Theta(\varepsilon^{-1})\) (Cybenko, 1989) or \(\Theta(\varepsilon^{-1/k})\) (Mhaskar, 1993); see Table~\ref{tab:comparison}.

\begin{table}[t]
\centering
\caption{Comparison of parameter count, working-precision bit complexity, and weight magnitudes for prior MLP constructions on analytic targets.}
\label{tab:comparison}
\small
\setlength{\tabcolsep}{4pt}
\resizebox{\columnwidth}{!}{%
\begin{tabular}{@{}lccc@{}}
\toprule
Method & \#params vs.\ \(\varepsilon\) & Bit complexity (working precision) & Weight magnitudes \\
\midrule
Cybenko (1989)
& \(\Theta(\varepsilon^{-1})\)
& \(\Theta(\log(1/\varepsilon))\)
& \(\|W\|_{\max} = \Theta\!\big(\varepsilon^{-1}\log(1/\varepsilon)\big)\) \\

Mhaskar (1993)
& \(\widetilde{O}\!\big(\log(1/\varepsilon)\log\log(1/\varepsilon)\big)\)
& \(\gtrsim \widetilde{\Theta}\!\big(\log(1/\varepsilon)\log\log(1/\varepsilon)\big)\)
& \(\|W\|_{\max} = \Theta(\varepsilon^{-1/k})\) (order \(k\)) \\

Quasi-interpolant (ours)
& \(O(\log(1/\varepsilon))\)
& \(O(\log(1/\varepsilon)+\log\log(1/\varepsilon))\)
& \(\|W\|_{\max} = O(\log(1/\varepsilon))\) \\
\bottomrule
\end{tabular}%
}
\end{table}

\paragraph{Predictions for trained networks.}
The construction makes three testable predictions about what a trained MLP \emph{should} look like to achieve high precision: (i) \(\lambda = \gamma h\) should plateau as width increases, requiring \(\gamma\) to grow proportionally with \(N\); (ii) the outer-layer weight magnitudes should remain \(O(1)\), not grow; and (iii) the learned features should distribute approximation power uniformly across neurons rather than concentrating it in a compressible subset. In Section~4, we verify that end-to-end trained MLPs violate all three: \(\gamma\) stays \(O(1)\) (driving \(\lambda \to 0\)), outer weights are too large, and learned representations exhibit rank saturation.

\subsection{Algorithm}
\label{subsec:qi-algorithm}

\begin{algorithm}[t]
\caption{Explicit Quasi-Interpolant MLP Construction}
\label{alg:qi-mlp}
\small
\begin{algorithmic}[1]
\Statex \textbf{Input:} grid size \(N\), bandwidth \(\gamma\), halo \(R\), stencil half-width \(K_c\), DFT size \(M\), activation \(\psi\) with \(K=\psi^{(r)}\), and target samples on the extended grid.

\State Build the extended uniform grid on \([-1,1]\) together with the halo points needed for the finite operator.
\State Sample the target \(f\) on this extended grid.
\State Evaluate the periodized kernel spectrum \(D_h\) on a DFT grid and form the Fourier character \(\widetilde{C}_\lambda = h/\widetilde{D}_\lambda\).
\State Recover the cardinal coefficients \(c_j\) by inverse DFT, and retain the stencil coefficients with \(|j| \le K_c\).
\State Form the truncated cardinal function \(L_h^{(K_c)}(x) = \sum_{|j|\le K_c} c_j K_\gamma(x-jh)\).
\State Convolve the stencil coefficients \(c_j\) with the sampled target values to obtain the kernel-network coefficients \(a[m]\).
\State Convert the kernel network into MLP parameters by using hidden units \(\psi(\gamma(x-x_m))\) and outer weights \(w[m] = a[m]/\gamma^{r-1}\).
\State If \(r>1\), determine the degree-\((r-1)\) integration polynomial from boundary conditions.
\State Output the final approximant \(\widetilde f(x) = p_{r-1}(x) + \sum_{m=-R}^{N+R} w[m]\psi(\gamma(x-x_m))\).

\Statex \textbf{Output:} the explicit one-hidden-layer MLP \(\widetilde f\).

\Statex \textbf{Width:} \(W = N + 2R + 1\) neurons, plus \(r\) polynomial coefficients when \(r>1\).
\end{algorithmic}
\end{algorithm}

% =========================
% PASTE MY SECTION 3 BLOCK HERE
% =========================

\section{Experiments}
% Paste your experiments here

\section{Discussion}
% Paste your discussion here

\medskip
\small

\bibliographystyle{plainnat}
\bibliography{references}

\appendix

\section{Theoretical Results}
% Keep your existing Appendix A here

% =========================
% PASTE MY APPENDIX B AND C BLOCKS HERE
% =========================

% The NeurIPS checklist is included by the style file.
% Leave it in place for submission.

% =========================
% APPENDICES B AND C ONLY
% Paste these after existing Appendix A.
% Assumes \appendix has already been called.
% =========================

\section{Construction Derivations}
\label{app:construction-derivations}

This appendix derives the Fourier normalization construction summarized in Section~3.1. The proof of Theorem~1 (with explicit constants), Lemmas~2--4, and the kernel verification for \(\operatorname{sech}^2\) are in Appendix~A.1.

\subsection{Deriving the Cardinal Function via Fourier Normalization}
\label{app:fourier-normalization}

We establish that the cardinal property \(L_h(jh) = \delta_{j,0}\) is equivalent to \eqref{eq:partition-unity}, and show how Fourier normalization yields equation~\eqref{eq:fourier-character}.

\subsubsection{The cardinal property in Fourier space}

Write
\[
L_h(jh) = \frac{1}{2\pi} \int_{\mathbb{R}} \widehat{L}_h(\omega)\, e^{i\omega jh}\, d\omega.
\]
Partition \(\mathbb{R}\) into intervals of length \(2\pi/h\) and substitute \(\omega \mapsto \omega + 2\pi m/h\) in each:
\[
L_h(jh)
=
\frac{1}{2\pi}
\int_{-\pi/h}^{\pi/h}
\underbrace{
\left[
\sum_{m \in \mathbb{Z}}
\widehat{L}_h\!\left(\omega + \frac{2\pi m}{h}\right)
\right]
}_{=: \, S(\omega)}
e^{i\omega jh}\, d\omega,
\]
using \(e^{i(2\pi m/h)jh} = 1\). The integral extracts the \(j\)-th Fourier coefficient of \(S\) on \([-\pi/h,\pi/h]\).

\paragraph{Forward.}
If \(S(\omega) = h\), then
\[
L_h(jh)
=
\frac{h}{2\pi}\int_{-\pi/h}^{\pi/h} e^{i\omega jh}\, d\omega
=
\delta_{j,0}
\]
by orthogonality.

\paragraph{Reverse.}
If \(L_h(jh) = \delta_{j,0}\) for all \(j\), all nonzero Fourier coefficients of \(S\) vanish, forcing \(S = h\).

\subsubsection{From kernel composition to equation~\eqref{eq:fourier-character}}

Write \(L_h(x) = \sum_j c_j K_\gamma(x - jh)\). Fourier-transforming gives
\[
\widehat{L}_h(\omega) = \widehat{K}_\gamma(\omega)\,\widehat{C}_h(\omega),
\qquad
\widehat{C}_h(\omega) := \sum_j c_j e^{-ijh\omega}.
\]
Because \(\widehat{C}_h\) is \(2\pi/h\)-periodic,
\[
\sum_{m \in \mathbb{Z}} \widehat{L}_h\!\left(\omega+\frac{2\pi m}{h}\right)
=
\widehat{C}_h(\omega)\sum_{m \in \mathbb{Z}} \widehat{K}_\gamma\!\left(\omega+\frac{2\pi m}{h}\right)
=
\widehat{C}_h(\omega)\,D_h(\omega).
\]
Plugging this into \eqref{eq:partition-unity} yields
\[
h = \widehat{C}_h(\omega)\,D_h(\omega),
\]
so
\[
\widehat{C}_h(\omega) = \frac{h}{D_h(\omega)}.
\]

\subsubsection{Recovering the coefficients \(c_j\)}

Rescale \(\theta := h\omega\), and define
\[
\widetilde{C}_\lambda(\theta)
:=
\widehat{C}_h(\theta/h)
=
\frac{h}{\widetilde{D}_\lambda(\theta)}.
\]
This is \(2\pi\)-periodic with Fourier series
\[
\widetilde{C}_\lambda(\theta) = \sum_j c_j e^{-ij\theta}.
\]
The Fourier coefficients of \(\widetilde{C}_\lambda\) are the kernel-expansion weights \(c_j\).

\paragraph{Remark.}
The rescaled normalizer \(\widetilde{D}_\lambda(\theta) = D_h(\theta/h)\) depends on \(\gamma\) and \(h\) only through \(\lambda = \gamma h\), which is why \(\lambda\) is the natural dimensionless parameter.

\subsection{DFT Approximation of the Cardinal Coefficients}
\label{app:dft-approximation}

The coefficients \(c_j\) are computed via the \(M\)-point DFT of \(\widetilde{C}_\lambda\) at nodes \(\theta_l = 2\pi l/M\) (Algorithm~\ref{alg:qi-mlp}, Step~2). The DFT returns aliased coefficients
\[
\widetilde{c}_j^{(M)} = c_j + \sum_{l \neq 0} c_{j+lM}.
\]
By Lemma~2, \(|c_j| \leq C_c(\lambda)e^{-\sigma|j|}\), so
\[
\bigl|\widetilde{c}_j^{(M)} - c_j\bigr|
\leq
\frac{2C_c(\lambda)\, e^{-\sigma M}}{1 - e^{-\sigma M}}.
\]
This computational error is independent of the four approximation-theoretic terms in Theorem~1 and can be driven below machine precision by choosing \(M\) moderately larger than \(2K_c + 1\).

\subsection{Kernel Network Re-indexing}
\label{app:kernel-reindexing}

Substituting \eqref{eq:Lh-trunc} into \eqref{eq:finite-op} and setting \(m = k + j\) gives
\[
(Q_{h,R}^{(K_c)} f)(x)
=
\sum_{k=-R}^{N+R} f(x_k)\sum_{|j| \le K_c} c_j K_\gamma(x - x_k - jh)
=
\sum_{m=-R}^{N+R} a[m]\, K_\gamma(x - x_m),
\]
where
\[
a[m] = \sum_{|j| \le K_c} c_j f(x_{m-j})
\]
is a discrete convolution of width \(2K_c + 1\). This requires samples on the extended index set \(\{-R-K_c, \ldots, N+R+K_c\}\).

\subsection{Activation Coverage}
\label{app:activation-coverage}

The derivative-kernel relation \(K = \psi^{(r)}\) is verified for \(\tanh\):

\begin{center}
\small
\begin{tabular}{@{}llll@{}}
\toprule
Activation \(\psi\) & Kernel \(K\) & Order \(r\) & Integration polynomial \\
\midrule
\(\tanh\) & \(\operatorname{sech}^2\) & \(1\) & bias \\
\bottomrule
\end{tabular}
\end{center}

For \(\tanh\), the kernel \(K = \operatorname{sech}^2\) satisfies all conditions (K1--K4) as verified in Appendix~A.1.3: nonzero mass, exponential spatial decay with \((A_K, a_K) = (2,2)\), positive Fourier transform, and a diagonal-dominance condition ensuring a zero-free strip.

\paragraph{Remark on other activations.}
In principle, any smooth activation whose \(r\)-th derivative satisfies K1--K4 yields a valid construction. However, extending the verification beyond \(\tanh\) requires care. For instance, \(\mathrm{GELU}''(x) = \phi(x)(2 - x^2)\), which changes sign and has nonzero integral only after checking cancellations. \(\mathrm{Softplus}'(x) = \sigma(x)\), which does not decay as \(x \to +\infty\) and therefore violates K4 as stated. We leave a systematic treatment of activation coverage to future work; all experiments in this paper use \(\tanh\) (equivalently, \(\operatorname{sech}^2\)) activations.

\section{Width-Precision Scaling}
\label{app:width-precision}

\subsection{The Viable Regime for \(\lambda\)}
\label{app:viable-lambda}

For a fixed target tolerance \(\varepsilon\), the bandwidth \(\lambda\) must satisfy competing constraints.

\paragraph{Aliasing (upper bound).}
\(C_2 e^{-c_2/\lambda^p} \leq \varepsilon\) requires
\[
\lambda \leq \lambda_{\max}(\varepsilon)
:=
\left(\frac{c_2}{\log(C_2/\varepsilon)}\right)^{1/p}.
\]

\paragraph{Conditioning (lower bound).}
The prefactors \(C_i(\lambda)\) diverge as \(\lambda \to 0\). For \(K = \tfrac{1}{2}\operatorname{sech}^2\), the cardinal function's spatial decay (Lemma~4) requires \(a_K \lambda > \sigma\), giving \(\lambda > \sigma/2\) with \(a_K = 2\). The diagonal-dominance condition (Proposition~4) further requires \(q(\sigma,\lambda) < 1\), satisfied for \(\lambda \gtrsim 1.2\) with \(\sigma = 1\).

The viable regime \([\lambda_{\min}, \lambda_{\max}(\varepsilon)]\) is nonempty for all \(\varepsilon\) above a kernel-dependent threshold; this is the threshold referenced in the informal main result. The threshold arises because \(\lambda_{\max}(\varepsilon) \to 0\) as \(\varepsilon \to 0\) while \(\lambda_{\min}\) is fixed by conditioning, so there exists a smallest achievable \(\varepsilon\) for each kernel. Note that \(\lambda_{\max}(\varepsilon)\) depends on \(\varepsilon\) only through a \(1/\log(1/\varepsilon)\) factor, so the viable regime shrinks slowly as \(\varepsilon \to 0\). For fp64 targets (\(\varepsilon \sim 10^{-15}\)) with the \(\operatorname{sech}^2\) kernel, the interval is well-separated and \(\lambda^* \approx 1.5\) is near-optimal across target functions in our experiments.

\paragraph{Practical selection.}
We choose \(\lambda^*\) by a one-dimensional sweep over \([1,4]\), evaluating the total bound \eqref{eq:error-bound} with all prefactors in closed form (Appendix~A.1, Eqs.~29--31).

\subsection{Proof of Corollary 1}
\label{app:proof-cor1}

Fix \(\varepsilon\) and choose \(\lambda^*\) in the viable regime so that the aliasing term satisfies
\[
C_2 e^{-c_2/(\lambda^*)^p} =: \delta \leq \varepsilon.
\]
Given width \(W\), allocate
\[
N = \lfloor \beta W \rfloor,
\qquad
R = \left\lfloor \frac{(1-\beta)W}{2} \right\rfloor,
\qquad
K_c = R
\]
for fixed \(\beta \in (0,1)\). Then \(h = 2/N \le 4/(\beta W)\), and:
\begin{itemize}
    \item \emph{Resolution:}
    \[
    C_1(\lambda^*) e^{-c_1 N/2}
    \le
    C_1(\lambda^*) e^{-c_1 \beta W/4}.
    \]

    \item \emph{Halo:}
    \[
    C_3(\lambda^*) e^{-c_3 \lambda^* R}
    \le
    C_3(\lambda^*) e^{-c_3 \lambda^*(1-\beta)W/4}.
    \]

    \item \emph{Stencil:}
    \[
    C_4(\lambda^*) e^{-c_4 \lambda^* R}
    \le
    C_4(\lambda^*) e^{-c_4 \lambda^*(1-\beta)W/4}.
    \]
\end{itemize}

Setting
\[
\alpha := \min\!\left(
\frac{c_1\beta}{4},
\frac{c_3\lambda^*(1-\beta)}{4},
\frac{c_4\lambda^*(1-\beta)}{4}
\right)
\]
and
\[
A := C_1 + C_3 + C_4 + \delta
\]
yields
\[
\|f - \widetilde{f}_W\| \leq A e^{-\alpha W}
\]
for \(W\) large enough that \(N \ge 2\) and \(R \ge 1\). \hfill\(\square\)

\paragraph{Remark.}
The constant \(A\) depends on \(\lambda^*\) and hence (weakly) on \(\varepsilon\). For any fixed precision target, \(A\) is a computable constant. The proportion \(\beta\) can be optimized:
\[
\beta^* = \frac{c_3\lambda^*}{c_1 + c_3\lambda^*};
\]
for typical parameters, \(\beta^* \approx 0.6\).

\subsection{Weight Magnitude Scaling}
\label{app:weight-scaling}

With \(\lambda^*\) fixed,
\[
\gamma = \lambda^* N/2.
\]

\paragraph{Inner-layer weights.}
\(\gamma = O(W) = O(\log(1/\varepsilon))\) by Corollary~1.

\paragraph{Outer-layer weights.}
\[
|a[m]|
\le
\|f\|_\infty \sum_{|j| \le K_c} |c_j|
\le
\|f\|_\infty C_c(\lambda^*) \coth(\sigma/2)
=
O(1)
\]
by Lemma~2. So \(w[m] = a[m]/\gamma^{r-1}\): for \(r = 1\) (\(\tanh\)), \(w[m] = O(1)\).

\paragraph{Maximum weight.}
\[
\|W\|_{\max} = \gamma = O(\log(1/\varepsilon)).
\]
Concretely, for \(\varepsilon = 10^{-15}\): our construction gives \(\gamma \sim 35\), versus \(\sim 10^{15}\) for Cybenko (1989) and \(\sim 10^{15/k}\) for Mhaskar (1993).

\end{document}